% ================================================================================
% MANUAL DE USUARIO — GAMMA
% Dirigido al gestor de datos maestros. No cubre funciones de administracion.
%
% Compilar con:  latexmk -pdf manual_usuario.tex
%
% CAPTURAS: van en capturas/ con los nombres que indica cada \captura.
% Mientras falte una, se dibuja un recuadro con el nombre esperado, de modo que
% el documento compila desde el primer dia y las imagenes se van agregando.
% ================================================================================
% ================================================================================
% PREAMBULO COMPARTIDO — Manuales de GAMMA
%
% Lo usan los tres manuales (usuario, tecnico e instalacion) para que compartan
% tipografia, paleta y recuadros. Se incluye con % ================================================================================
% PREAMBULO COMPARTIDO — Manuales de GAMMA
%
% Lo usan los tres manuales (usuario, tecnico e instalacion) para que compartan
% tipografia, paleta y recuadros. Se incluye con % ================================================================================
% PREAMBULO COMPARTIDO — Manuales de GAMMA
%
% Lo usan los tres manuales (usuario, tecnico e instalacion) para que compartan
% tipografia, paleta y recuadros. Se incluye con \input{../preambulo} desde la
% carpeta de cada manual.
%
% La paleta sale de las variables de frontend/src/style.css: la interfaz es
% enteramente acromatica, y los manuales tampoco introducen color.
% ================================================================================
\documentclass[11pt, letterpaper]{article}

\usepackage[spanish,mexico]{babel}
\usepackage[utf8]{inputenc}
\usepackage[T1]{fontenc}

% Tipografia: Inter, la mas cercana en TeX Live a la Geist que usa la
% aplicacion. Sin serifas en todo el documento, igual que la interfaz.
\usepackage[sfdefault]{inter}
\usepackage[scaled=0.86]{beramono}

\usepackage[top=2.4cm, bottom=2.4cm, left=2.6cm, right=2.6cm]{geometry}
\usepackage{graphicx}
\usepackage{xcolor}
\usepackage{enumitem}
\usepackage{fancyhdr}
\usepackage{titlesec}
\usepackage{tcolorbox}
% `enhanced` y `borderline west` viven en skins; `breakable` en su propia
% libreria. Sin declararlas, tcolorbox ignora esas claves en silencio y los
% recuadros pierden la barra lateral que los distingue.
\tcbuselibrary{skins,breakable}
\usepackage{booktabs}
\usepackage[hidelinks]{hyperref}

% ── Paleta ──────────────────────────────────────────────────────────────────
% Tomada de las variables de src/style.css. La interfaz es enteramente
% acromatica —todos sus colores son oklch(L 0 0)— asi que el manual tampoco
% introduce color: solo grises, con el texto como unico elemento de contraste.
\definecolor{fg}{HTML}{0A0A0A}        % --foreground
\definecolor{fgsoft}{HTML}{171717}    % --primary
\definecolor{muted}{HTML}{737373}     % --muted-foreground
\definecolor{surface}{HTML}{F5F5F5}   % --muted / --secondary
\definecolor{surfacealt}{HTML}{FAFAFA}% --sidebar
\definecolor{line}{HTML}{E5E5E5}      % --border

\titleformat{\section}{\LARGE\bfseries\color{fg}}{\thesection}{0.7em}{}
\titlespacing*{\section}{0pt}{1.6em}{0.7em}
\titleformat{\subsection}{\large\bfseries\color{fgsoft}}{\thesubsection}{0.6em}{}
\titlespacing*{\subsection}{0pt}{1.3em}{0.4em}
\titleformat{\subsubsection}[runin]{\bfseries\color{fgsoft}}{}{0pt}{}[.]

\pagestyle{fancy}
\fancyhf{}
\lhead{\footnotesize\color{muted}Manual de usuario}
\rhead{\footnotesize\color{muted}\thepage}
\renewcommand{\headrulewidth}{0.4pt}
\renewcommand{\headrule}{\color{line}\hrule height 0.4pt}

\setlength{\parindent}{0pt}
\setlength{\parskip}{0.62em}

% Enlaces y referencias en el mismo gris del cuerpo: sin subrayados ni azules.
\hypersetup{colorlinks=false}

% ── Recuadros ───────────────────────────────────────────────────────────────
% Se distinguen por una barra lateral y el peso del titulo, no por color.
\newtcolorbox{nota}[1][Nota]{
  enhanced, breakable, colback=surfacealt, colframe=surfacealt,
  borderline west={2pt}{0pt}{line},
  coltitle=fg, fonttitle=\bfseries\footnotesize\sffamily,
  title=#1, boxrule=0pt, arc=0pt, outer arc=0pt,
  left=10pt, right=10pt, top=7pt, bottom=7pt, toptitle=1pt, bottomtitle=3pt}

\newtcolorbox{aviso}[1][Importante]{
  enhanced, breakable, colback=surface, colframe=surface,
  borderline west={2pt}{0pt}{fg},
  coltitle=fg, fonttitle=\bfseries\footnotesize\sffamily,
  title=#1, boxrule=0pt, arc=0pt, outer arc=0pt,
  left=10pt, right=10pt, top=7pt, bottom=7pt, toptitle=1pt, bottomtitle=3pt}

% ── Captura de pantalla con marcador de posicion ────────────────────────────
% Uso: \captura{nombre-archivo}{Pie de la figura}
% Si manual/capturas/<nombre>.png existe, la inserta. Si no, dibuja el hueco.
\newcommand{\captura}[2]{%
  \begin{figure}[htbp]
  \centering
  \IfFileExists{capturas/#1.png}{%
    % Ancho completo, pero acotado en alto: una captura de ventana completa
    % puede ser muy alta y desbordaria la pagina. keepaspectratio evita que se
    % deforme cuando manda la restriccion de altura.
    {\color{line}\fboxrule=0.6pt\fboxsep=0pt\fbox{%
      \includegraphics[width=0.94\textwidth, height=0.62\textheight, keepaspectratio]{capturas/#1.png}}}%
  }{%
    \fcolorbox{line}{surfacealt}{%
      \begin{minipage}[c][3.1cm][c]{0.93\textwidth}
        \centering
        {\footnotesize\bfseries\color{muted}CAPTURA PENDIENTE}\\[5pt]
        {\ttfamily\footnotesize\color{fgsoft} capturas/#1.png}\\[4pt]
        {\footnotesize\color{muted} #2}
      \end{minipage}}%
  }
  \caption{#2}
  \label{fig:#1}
  \end{figure}%
}

% Pies de figura discretos
\usepackage{caption}
\captionsetup{font={small,color=muted}, labelfont={bf,color=muted}, skip=6pt}

 desde la
% carpeta de cada manual.
%
% La paleta sale de las variables de frontend/src/style.css: la interfaz es
% enteramente acromatica, y los manuales tampoco introducen color.
% ================================================================================
\documentclass[11pt, letterpaper]{article}

\usepackage[spanish,mexico]{babel}
\usepackage[utf8]{inputenc}
\usepackage[T1]{fontenc}

% Tipografia: Inter, la mas cercana en TeX Live a la Geist que usa la
% aplicacion. Sin serifas en todo el documento, igual que la interfaz.
\usepackage[sfdefault]{inter}
\usepackage[scaled=0.86]{beramono}

\usepackage[top=2.4cm, bottom=2.4cm, left=2.6cm, right=2.6cm]{geometry}
\usepackage{graphicx}
\usepackage{xcolor}
\usepackage{enumitem}
\usepackage{fancyhdr}
\usepackage{titlesec}
\usepackage{tcolorbox}
% `enhanced` y `borderline west` viven en skins; `breakable` en su propia
% libreria. Sin declararlas, tcolorbox ignora esas claves en silencio y los
% recuadros pierden la barra lateral que los distingue.
\tcbuselibrary{skins,breakable}
\usepackage{booktabs}
\usepackage[hidelinks]{hyperref}

% ── Paleta ──────────────────────────────────────────────────────────────────
% Tomada de las variables de src/style.css. La interfaz es enteramente
% acromatica —todos sus colores son oklch(L 0 0)— asi que el manual tampoco
% introduce color: solo grises, con el texto como unico elemento de contraste.
\definecolor{fg}{HTML}{0A0A0A}        % --foreground
\definecolor{fgsoft}{HTML}{171717}    % --primary
\definecolor{muted}{HTML}{737373}     % --muted-foreground
\definecolor{surface}{HTML}{F5F5F5}   % --muted / --secondary
\definecolor{surfacealt}{HTML}{FAFAFA}% --sidebar
\definecolor{line}{HTML}{E5E5E5}      % --border

\titleformat{\section}{\LARGE\bfseries\color{fg}}{\thesection}{0.7em}{}
\titlespacing*{\section}{0pt}{1.6em}{0.7em}
\titleformat{\subsection}{\large\bfseries\color{fgsoft}}{\thesubsection}{0.6em}{}
\titlespacing*{\subsection}{0pt}{1.3em}{0.4em}
\titleformat{\subsubsection}[runin]{\bfseries\color{fgsoft}}{}{0pt}{}[.]

\pagestyle{fancy}
\fancyhf{}
\lhead{\footnotesize\color{muted}Manual de usuario}
\rhead{\footnotesize\color{muted}\thepage}
\renewcommand{\headrulewidth}{0.4pt}
\renewcommand{\headrule}{\color{line}\hrule height 0.4pt}

\setlength{\parindent}{0pt}
\setlength{\parskip}{0.62em}

% Enlaces y referencias en el mismo gris del cuerpo: sin subrayados ni azules.
\hypersetup{colorlinks=false}

% ── Recuadros ───────────────────────────────────────────────────────────────
% Se distinguen por una barra lateral y el peso del titulo, no por color.
\newtcolorbox{nota}[1][Nota]{
  enhanced, breakable, colback=surfacealt, colframe=surfacealt,
  borderline west={2pt}{0pt}{line},
  coltitle=fg, fonttitle=\bfseries\footnotesize\sffamily,
  title=#1, boxrule=0pt, arc=0pt, outer arc=0pt,
  left=10pt, right=10pt, top=7pt, bottom=7pt, toptitle=1pt, bottomtitle=3pt}

\newtcolorbox{aviso}[1][Importante]{
  enhanced, breakable, colback=surface, colframe=surface,
  borderline west={2pt}{0pt}{fg},
  coltitle=fg, fonttitle=\bfseries\footnotesize\sffamily,
  title=#1, boxrule=0pt, arc=0pt, outer arc=0pt,
  left=10pt, right=10pt, top=7pt, bottom=7pt, toptitle=1pt, bottomtitle=3pt}

% ── Captura de pantalla con marcador de posicion ────────────────────────────
% Uso: \captura{nombre-archivo}{Pie de la figura}
% Si manual/capturas/<nombre>.png existe, la inserta. Si no, dibuja el hueco.
\newcommand{\captura}[2]{%
  \begin{figure}[htbp]
  \centering
  \IfFileExists{capturas/#1.png}{%
    % Ancho completo, pero acotado en alto: una captura de ventana completa
    % puede ser muy alta y desbordaria la pagina. keepaspectratio evita que se
    % deforme cuando manda la restriccion de altura.
    {\color{line}\fboxrule=0.6pt\fboxsep=0pt\fbox{%
      \includegraphics[width=0.94\textwidth, height=0.62\textheight, keepaspectratio]{capturas/#1.png}}}%
  }{%
    \fcolorbox{line}{surfacealt}{%
      \begin{minipage}[c][3.1cm][c]{0.93\textwidth}
        \centering
        {\footnotesize\bfseries\color{muted}CAPTURA PENDIENTE}\\[5pt]
        {\ttfamily\footnotesize\color{fgsoft} capturas/#1.png}\\[4pt]
        {\footnotesize\color{muted} #2}
      \end{minipage}}%
  }
  \caption{#2}
  \label{fig:#1}
  \end{figure}%
}

% Pies de figura discretos
\usepackage{caption}
\captionsetup{font={small,color=muted}, labelfont={bf,color=muted}, skip=6pt}

 desde la
% carpeta de cada manual.
%
% La paleta sale de las variables de frontend/src/style.css: la interfaz es
% enteramente acromatica, y los manuales tampoco introducen color.
% ================================================================================
\documentclass[11pt, letterpaper]{article}

\usepackage[spanish,mexico]{babel}
\usepackage[utf8]{inputenc}
\usepackage[T1]{fontenc}

% Tipografia: Inter, la mas cercana en TeX Live a la Geist que usa la
% aplicacion. Sin serifas en todo el documento, igual que la interfaz.
\usepackage[sfdefault]{inter}
\usepackage[scaled=0.86]{beramono}

\usepackage[top=2.4cm, bottom=2.4cm, left=2.6cm, right=2.6cm]{geometry}
\usepackage{graphicx}
\usepackage{xcolor}
\usepackage{enumitem}
\usepackage{fancyhdr}
\usepackage{titlesec}
\usepackage{tcolorbox}
% `enhanced` y `borderline west` viven en skins; `breakable` en su propia
% libreria. Sin declararlas, tcolorbox ignora esas claves en silencio y los
% recuadros pierden la barra lateral que los distingue.
\tcbuselibrary{skins,breakable}
\usepackage{booktabs}
\usepackage[hidelinks]{hyperref}

% ── Paleta ──────────────────────────────────────────────────────────────────
% Tomada de las variables de src/style.css. La interfaz es enteramente
% acromatica —todos sus colores son oklch(L 0 0)— asi que el manual tampoco
% introduce color: solo grises, con el texto como unico elemento de contraste.
\definecolor{fg}{HTML}{0A0A0A}        % --foreground
\definecolor{fgsoft}{HTML}{171717}    % --primary
\definecolor{muted}{HTML}{737373}     % --muted-foreground
\definecolor{surface}{HTML}{F5F5F5}   % --muted / --secondary
\definecolor{surfacealt}{HTML}{FAFAFA}% --sidebar
\definecolor{line}{HTML}{E5E5E5}      % --border

\titleformat{\section}{\LARGE\bfseries\color{fg}}{\thesection}{0.7em}{}
\titlespacing*{\section}{0pt}{1.6em}{0.7em}
\titleformat{\subsection}{\large\bfseries\color{fgsoft}}{\thesubsection}{0.6em}{}
\titlespacing*{\subsection}{0pt}{1.3em}{0.4em}
\titleformat{\subsubsection}[runin]{\bfseries\color{fgsoft}}{}{0pt}{}[.]

\pagestyle{fancy}
\fancyhf{}
\lhead{\footnotesize\color{muted}Manual de usuario}
\rhead{\footnotesize\color{muted}\thepage}
\renewcommand{\headrulewidth}{0.4pt}
\renewcommand{\headrule}{\color{line}\hrule height 0.4pt}

\setlength{\parindent}{0pt}
\setlength{\parskip}{0.62em}

% Enlaces y referencias en el mismo gris del cuerpo: sin subrayados ni azules.
\hypersetup{colorlinks=false}

% ── Recuadros ───────────────────────────────────────────────────────────────
% Se distinguen por una barra lateral y el peso del titulo, no por color.
\newtcolorbox{nota}[1][Nota]{
  enhanced, breakable, colback=surfacealt, colframe=surfacealt,
  borderline west={2pt}{0pt}{line},
  coltitle=fg, fonttitle=\bfseries\footnotesize\sffamily,
  title=#1, boxrule=0pt, arc=0pt, outer arc=0pt,
  left=10pt, right=10pt, top=7pt, bottom=7pt, toptitle=1pt, bottomtitle=3pt}

\newtcolorbox{aviso}[1][Importante]{
  enhanced, breakable, colback=surface, colframe=surface,
  borderline west={2pt}{0pt}{fg},
  coltitle=fg, fonttitle=\bfseries\footnotesize\sffamily,
  title=#1, boxrule=0pt, arc=0pt, outer arc=0pt,
  left=10pt, right=10pt, top=7pt, bottom=7pt, toptitle=1pt, bottomtitle=3pt}

% ── Captura de pantalla con marcador de posicion ────────────────────────────
% Uso: \captura{nombre-archivo}{Pie de la figura}
% Si manual/capturas/<nombre>.png existe, la inserta. Si no, dibuja el hueco.
\newcommand{\captura}[2]{%
  \begin{figure}[htbp]
  \centering
  \IfFileExists{capturas/#1.png}{%
    % Ancho completo, pero acotado en alto: una captura de ventana completa
    % puede ser muy alta y desbordaria la pagina. keepaspectratio evita que se
    % deforme cuando manda la restriccion de altura.
    {\color{line}\fboxrule=0.6pt\fboxsep=0pt\fbox{%
      \includegraphics[width=0.94\textwidth, height=0.62\textheight, keepaspectratio]{capturas/#1.png}}}%
  }{%
    \fcolorbox{line}{surfacealt}{%
      \begin{minipage}[c][3.1cm][c]{0.93\textwidth}
        \centering
        {\footnotesize\bfseries\color{muted}CAPTURA PENDIENTE}\\[5pt]
        {\ttfamily\footnotesize\color{fgsoft} capturas/#1.png}\\[4pt]
        {\footnotesize\color{muted} #2}
      \end{minipage}}%
  }
  \caption{#2}
  \label{fig:#1}
  \end{figure}%
}

% Pies de figura discretos
\usepackage{caption}
\captionsetup{font={small,color=muted}, labelfont={bf,color=muted}, skip=6pt}



\begin{document}

% ── Portada ─────────────────────────────────────────────────────────────────
\begin{titlepage}
\thispagestyle{empty}
\vspace*{5.5cm}
{\fontsize{46}{50}\selectfont\bfseries\color{fg} GAMMA\par}
\vspace{0.5em}
{\large\color{muted} Gobierno Automatizado del Maestro de Materiales\par}

\vspace{2.6cm}
{\color{line}\hrule height 1pt}
\vspace{1.2em}

{\huge\bfseries\color{fg} Manual de usuario\par}
\vspace{0.5em}
{\large\color{muted} Guía para el gestor de datos maestros\par}

\vfill
{\color{line}\hrule height 0.4pt}
\vspace{1em}
{\footnotesize\color{muted}
Unidad de Energía \quad·\quad \today\par
\vspace{0.4em}
Este manual cubre el uso diario de la herramienta.
Las funciones de administración se documentan por separado.\par}
\vspace{1cm}
\end{titlepage}

\tableofcontents
\newpage

% ============================================================================
\section{Qué es GAMMA y qué no es}

GAMMA es un asistente que lo acompaña al dar de alta un material en SAP. Hace
tres cosas por usted:

\begin{enumerate}[leftmargin=*, itemsep=3pt]
  \item \textbf{Redacta la descripción} en el formato estándar de la
        corporación, a partir de lo que usted escriba en lenguaje corriente.
  \item \textbf{Revisa si el material ya existe} comparándolo contra el
        catálogo, para evitar que se creen duplicados.
  \item \textbf{Sugiere la clase} a la que pertenece el material, con base en
        cómo se han clasificado materiales parecidos en el pasado.
\end{enumerate}

\begin{aviso}[GAMMA no crea el material en SAP]
La herramienta \textbf{no escribe en SAP}. Prepara y valida la solicitud; el
registro final lo hace usted, cargando en SAP el archivo de Excel que GAMMA
genera. Usted conserva siempre la última palabra: puede aceptar, corregir o
descartar cualquier sugerencia.
\end{aviso}

\subsection{El recorrido completo}

\begin{enumerate}[leftmargin=*, itemsep=3pt]
  \item Inicia sesión.
  \item Le describe el material al asistente, con sus palabras.
  \item Si falta información, el asistente le indica exactamente qué falta.
  \item Revisa los posibles duplicados que le muestre.
  \item Confirma o corrige la clase sugerida.
  \item Confirma la solicitud.
  \item Cuando tenga varias listas, las exporta a Excel y las carga en SAP.
\end{enumerate}

\newpage
% ============================================================================
\section{Iniciar sesión}

Abra la dirección de GAMMA en su navegador. Verá la pantalla de acceso.

\captura{01-login}{Pantalla de inicio de sesión}

Escriba su correo y su contraseña, y presione \textbf{Iniciar sesión}.

\begin{nota}[Si no tiene cuenta]
Las cuentas las crea el área de datos maestros. Si es la primera vez que entra
o olvidó su contraseña, solicítelo a su coordinación; el asistente no puede
crear cuentas.
\end{nota}

Una vez dentro, la barra superior le da acceso a todo. De izquierda a derecha
encontrará \textbf{Datos}, el botón de \textbf{Chat}, y luego tres controles:

\begin{itemize}[leftmargin=*, itemsep=3pt]
  \item El signo de interrogación abre la \textbf{Ayuda}.
  \item El sol o la luna cambia entre \textbf{modo claro y modo oscuro}. Su
        preferencia se recuerda para la próxima vez que entre.
  \item Su nombre despliega el \textbf{menú de la cuenta}.
\end{itemize}

\captura{02-barra-superior}{Barra de navegación, con la sesión iniciada}

\begin{nota}[Si ve más secciones de las que aquí se describen]
Las cuentas administrativas tienen acceso a secciones técnicas adicionales
—Arquitectura, Referencia, Lab, Modelo y Usuarios— que no forman parte del
trabajo diario del gestor y no se cubren en este manual.
\end{nota}

\subsection{Cerrar sesión}

Haga clic sobre su nombre, en el extremo derecho de la barra. Se despliega un
menú con sus datos y la opción \textbf{Cerrar sesión}.

\captura{03-menu-cuenta}{Menú de la cuenta}

\newpage
% ============================================================================
\section{Dar de alta un material}

Esta es la función principal. Entre a \textbf{Chat} desde la barra superior.

\captura{04-chat-inicio}{Pantalla del asistente al abrirla}

A la izquierda está el listado de sus conversaciones anteriores. El botón
\textbf{Nueva conversación} abre una en blanco. Conviene usar una conversación
por material: así queda ordenado y puede volver a consultarla después.

\subsection{Paso 1. Describa el material}

Escriba en el cuadro de abajo lo que necesita dar de alta, con sus palabras. No
tiene que usar abreviaturas, mayúsculas ni el formato de SAP: de eso se encarga
el asistente.

\begin{nota}[Ejemplos de cómo pedirlo]
\ttfamily
necesito una valvula de bola de pvc de 3 pulgadas\\[3pt]
tornillo hexagonal de acero galvanizado 1/2 por 2 pulgadas\\[3pt]
filtro de aceite para motor john deere, numero AT35155
\end{nota}

Mientras el asistente trabaja verá el avance en pantalla: primero redacta la
descripción, luego busca duplicados y por último sugiere la clase.

\captura{05-chat-procesando}{Indicador de avance mientras el asistente procesa}

\subsection{Paso 2. Complete lo que falte}

Si lo que escribió no alcanza, el asistente \textbf{no inventa datos}: le
responde con la lista de lo que hace falta.

\captura{06-chat-datos-faltantes}{Respuesta del asistente cuando falta información}

La respuesta viene redactada para que pueda copiarla y enviársela tal cual al
solicitante. Cuando le respondan, escriba los datos en la misma conversación y
el asistente continúa.

\begin{nota}[Dos campos que casi siempre le van a pedir]
El \textbf{material de fabricación} (acero, PVC, cobre, aluminio\ldots) es
obligatorio para materiales físicos.

La \textbf{marca y el modelo} solo se piden cuando se trata de un repuesto o un
componente donde la marca es indispensable para identificarlo. Para
herramientas, ferretería y consumibles genéricos no se piden.
\end{nota}

\subsection{Paso 3. Revise los posibles duplicados}

Si el asistente encuentra materiales parecidos en el catálogo, se los muestra
antes de continuar, con un porcentaje de parecido al lado de cada uno.

\captura{07-chat-duplicados}{Listado de posibles duplicados}

Tiene dos caminos:

\begin{itemize}[leftmargin=*, itemsep=3pt]
  \item \textbf{Alguno le sirve.} Haga clic sobre él. La solicitud se cierra
        indicando que el material ya existe, y usted se lleva su código. No se
        crea nada nuevo.
  \item \textbf{Ninguno le sirve.} Presione \textbf{Ninguno, continuar}. El
        asistente sigue con la propuesta de material nuevo.
\end{itemize}

\begin{nota}[Si le muestra cosas que no se parecen]
Es normal y es a propósito. El sistema prefiere mostrarle una sugerencia de más
antes que dejar pasar un duplicado. Revisarlas toma segundos; un duplicado en
el catálogo queda para siempre.
\end{nota}

\subsection{Paso 4. Revise la propuesta y la clase}

El asistente le presenta la propuesta completa: la descripción corta en formato
SAP, la descripción larga, el tipo de material y la clase sugerida con su grupo
de artículo y sector.

\captura{08-chat-propuesta}{Tarjeta con la propuesta de material}

Junto a la clase verá un porcentaje de confianza. Es qué tan seguro está el
modelo de esa sugerencia.

\begin{itemize}[leftmargin=*, itemsep=3pt]
  \item \textbf{Confianza alta.} La sugerencia es muy probablemente correcta.
        Reviséla de todas formas.
  \item \textbf{Confianza baja.} Ponga más atención. El asistente le ofrece
        alternativas.
\end{itemize}

\subsubsection*{Si la clase sugerida no es la correcta}

Puede buscar otra. Use el buscador de clases: escriba parte del código o del
nombre y seleccione la que corresponda.

\captura{09-chat-buscar-clase}{Buscador de clases}

\begin{nota}[Sus correcciones sirven]
Cuando usted corrige una clase, el sistema lo registra. Esas correcciones se
usan después para mejorar el modelo. Vale la pena tomarse el momento de
corregir en lugar de aceptar algo que no cuadra.
\end{nota}

\subsection{Paso 5. Confirme o descarte}

Al pie de la propuesta hay dos botones:

\begin{itemize}[leftmargin=*, itemsep=3pt]
  \item \textbf{Confirmar} — la solicitud queda lista para exportar a SAP.
  \item \textbf{Descartar} — la solicitud no procede y se cierra sin crear nada.
\end{itemize}

\captura{10-chat-confirmado}{Solicitud confirmada}

Una vez confirmada, la solicitud entra a la lista de pendientes por exportar.

\newpage
% ============================================================================
\section{Exportar a SAP}

Las solicitudes confirmadas no viajan solas a SAP: usted las descarga en un
archivo de Excel y lo carga en el sistema.

Entre a \textbf{Datos} en la barra superior y elija la pestaña
\textbf{Exportar}.

\captura{11-datos-exportar}{Pestaña de exportación}

\subsection{Cómo exportar}

\begin{enumerate}[leftmargin=*, itemsep=3pt]
  \item Elija el estado que quiere ver. Por omisión aparecen las
        \textbf{confirmadas}.
  \item Deje marcada la casilla \textbf{Excluir ya exportados} para no repetir
        lo que ya bajó antes.
  \item Marque las solicitudes que quiera incluir. Puede seleccionarlas todas.
  \item Presione \textbf{Exportar}. El archivo se descarga a su computadora.
\end{enumerate}

\captura{12-datos-exportar-seleccion}{Selección de solicitudes para exportar}

El archivo trae la descripción corta, la descripción larga, el tipo de material
y la clase asignada de cada solicitud: todo lo que necesita para el alta en SAP.

\begin{aviso}[Después de cargar en SAP]
Las solicitudes exportadas quedan marcadas como tales. Si vuelve a exportar con
la casilla \textbf{Excluir ya exportados} marcada, no le aparecerán otra vez.
Eso evita cargas duplicadas.
\end{aviso}

\newpage
% ============================================================================
\section{Importar el maestro de materiales}

La pestaña \textbf{Importar} de la sección \textbf{Datos} sirve para actualizar
la copia del catálogo con la que trabaja GAMMA. Es una tarea que se hace de
tanto en tanto, no a diario, y normalmente la coordina el área de datos
maestros.

\captura{13-datos-importar}{Pestaña de importación}

\subsection{Los tres archivos y su orden}

Se cargan tres cosas distintas, y \textbf{el orden importa}:

\begin{enumerate}[leftmargin=*, itemsep=4pt]
  \item \textbf{UNSPSC.} El clasificador de Naciones Unidas. Va primero porque
        las clases se apoyan en él.
  \item \textbf{Clases o denominaciones.} El catálogo de clases de material.
        Necesita el UNSPSC ya cargado.
  \item \textbf{Maestro de materiales.} Los archivos exportados de SAP, uno por
        tipo de material. Necesitan las clases ya cargadas.
\end{enumerate}

\begin{aviso}[Respete el orden]
Si carga los materiales antes que las clases, los que apunten a una clase que
todavía no existe entran \textbf{sin clasificar}. No dan error: simplemente
quedan fuera del alcance del asistente hasta que se vuelvan a cargar. Cargue
siempre UNSPSC, luego clases y al final materiales.
\end{aviso}

\subsection{Cómo cargar}

\begin{enumerate}[leftmargin=*, itemsep=3pt]
  \item Presione \textbf{Seleccionar archivos .xlsx} en el bloque que
        corresponda. Puede elegir varios de una vez.
  \item Los archivos se van a una cola y se procesan uno por uno. Cada línea
        muestra su avance y si terminó bien o falló.
  \item Al terminar, revise los conteos de la parte superior: deben haber
        subido.
\end{enumerate}

\begin{nota}[Cargar dos veces no duplica]
Los materiales se identifican por su código. Si vuelve a cargar un archivo que
ya subió, los registros existentes se actualizan en lugar de duplicarse, así
que repetir una carga es inofensivo.
\end{nota}

\newpage
% ============================================================================
\section{La sección de Ayuda}

El signo de interrogación de la barra superior abre una versión resumida de
esta guía, siempre a la mano y sin salir de la aplicación.

\captura{14-ayuda}{Sección de ayuda}

Contiene tres cosas:

\begin{itemize}[leftmargin=*, itemsep=3pt]
  \item \textbf{El manual completo}, para descargar en PDF. Es este mismo
        documento.
  \item \textbf{Una guía rápida} de los cinco pasos del flujo de creación, para
        consultar de un vistazo.
  \item \textbf{Preguntas frecuentes}, las mismas que encontrará más adelante.
\end{itemize}

Al pie hay un acceso directo al asistente, por si llegó buscando ayuda y lo que
necesitaba era empezar un alta.

\newpage
% ============================================================================
\section{Preguntas frecuentes}

\subsection*{Busqué un material y no apareció, pero sé que existe en SAP}

GAMMA trabaja sobre una copia del catálogo que se actualiza cada cierto tiempo,
no sobre SAP en vivo. Un material creado en SAP después de la última
actualización todavía no se ve. Si tiene dudas, verifique directamente en SAP
antes de crear.

\subsection*{¿Puedo corregir una solicitud que ya confirmé?}

Sí, mientras no la haya exportado: vuelva a la conversación y ajústela. Si ya la
exportó y la cargó en SAP, la corrección se hace en SAP.

\subsection*{El asistente me pidió datos que el solicitante no tiene}

Copie la lista que le dio el asistente y devuélvasela al solicitante. Si un dato
de verdad no existe o no aplica, indíqueselo al asistente en la conversación y
explique por qué; con eso continúa.

\subsection*{¿Qué tan confiable es la clase que sugiere?}

Acierta en cerca de ocho de cada diez casos en su primera sugerencia, y la clase
correcta aparece entre las tres primeras en más de nueve de cada diez. Por eso
la decisión sigue siendo suya: el asistente propone, usted valida.

\subsection*{¿Se pierde mi conversación si cierro el navegador?}

No. Las conversaciones quedan guardadas y las encuentra en la lista de la
izquierda cuando vuelva a entrar.

\subsection*{¿Puedo trabajar varios materiales a la vez?}

Sí, pero abra una conversación por material. Mezclar varios en la misma
conversación confunde al asistente y complica el seguimiento.

\begin{nota}[¿Necesita ayuda?]
Si algo no funciona como esperaba, escriba al área de datos maestros
describiendo el material y el paso en el que se quedó. Si el asistente le
mostró un mensaje de error, cópielo tal cual: ayuda mucho a resolverlo.
\end{nota}

\end{document}
